% ============================================================
\chapter{Discrete Dynamical Systems}
\label{ch:discrete}
% ============================================================

\section{Introduction}

Discrete dynamical systems\index{discrete dynamical system}, defined by the
iteration of a map $x_{n+1} = f(x_n)$, provide the simplest framework for
studying chaos. This chapter explores the fundamental phenomena:
period-doubling cascades, universality, symbolic dynamics, and topological
entropy.

\begin{intuition}
Even a map as simple as $f(x) = rx(1-x)$ generates, as $r$ increases, the
full complexity of chaos: bifurcations, self-similarity, and universality.
Discrete systems are the ideal laboratory for understanding fundamental
mechanisms.
\end{intuition}

\section{The Logistic Map}

\begin{definition}[Logistic family]
\index{logistic map}
The \textbf{logistic family} is $f_r : [0,1] \to [0,1]$ defined by
$f_r(x) = rx(1-x)$, $r \in [0, 4]$.
\end{definition}

\begin{proposition}[Fixed points and stability]
\begin{enumerate}
    \item The fixed point $x^* = 0$ exists for all $r$; it is stable if $r < 1$.
    \item The fixed point $x^* = 1 - 1/r$ exists for $r > 1$; it is stable if
    $1 < r < 3$ (since $f_r'(x^*) = 2 - r$).
    \item At $r = 3$, a period-doubling bifurcation creates a period-2 cycle.
\end{enumerate}
\end{proposition}

\begin{theorem}[Period-doubling cascade]
\index{period-doubling cascade}
There exists an increasing sequence $r_1 < r_2 < r_3 < \cdots$ with $r_1 = 3$
such that at $r = r_n$, the period-$2^{n-1}$ cycle loses stability and a
stable period-$2^n$ cycle appears via a flip bifurcation. The sequence $r_n$
converges to $r_\infty \approx 3.5699\ldots$
\end{theorem}

\begin{theorem}[Feigenbaum universality]
\index{Feigenbaum constant}
The ratio $\delta_n = \frac{r_n - r_{n-1}}{r_{n+1} - r_n}$ converges to the
universal Feigenbaum constant
\[
    \delta_F = \lim_{n \to \infty} \delta_n = 4.669\,201\,609\ldots
\]
This constant is the same for every unimodal family with a quadratic maximum.
The spatial scaling factor converges to $\alpha_F = -2.502\,907\,875\ldots$
\end{theorem}

\begin{keyformulas}[title=\textbf{Regimes of the Logistic Map}]
\begin{align*}
    r \in [0, 1] &: \text{convergence to } 0. \\
    r \in (1, 3) &: \text{convergence to } x^* = 1 - 1/r. \\
    r \in (3, 1{+}\sqrt{6}) &: \text{stable period-2 cycle.} \\
    r = r_\infty \approx 3.5699 &: \text{accumulation of doublings.} \\
    r \in (r_\infty, 4) &: \text{mixture of chaos and periodic windows.} \\
    r = 4 &: \text{fully developed chaos on } [0, 1].
\end{align*}
\end{keyformulas}

\section{Sharkovskii's Theorem}

\begin{definition}[Sharkovskii ordering]
\index{Sharkovskii ordering}
The \textbf{Sharkovskii ordering} on $\N^*$ is:
\begin{multline*}
    3 \triangleright 5 \triangleright 7 \triangleright 9 \triangleright \cdots
    \triangleright 2{\cdot}3 \triangleright 2{\cdot}5 \triangleright 2{\cdot}7
    \triangleright \cdots \\
    \triangleright 4{\cdot}3 \triangleright 4{\cdot}5 \triangleright \cdots
    \triangleright \cdots \triangleright 2^3 \triangleright 2^2
    \triangleright 2 \triangleright 1.
\end{multline*}
\end{definition}

\begin{theorem}[Sharkovskii, 1964]
\index{Sharkovskii's theorem}
Let $f : \R \to \R$ be continuous. If $f$ has a periodic point of period $m$,
then $f$ has periodic points of every period $n$ with $m \triangleright n$ in
the Sharkovskii ordering. Moreover, for each $m$, there exists a continuous
map having a period-$m$ point but no period-$n$ point for $n \triangleright m$.
\end{theorem}

\begin{corollary}[Period three implies all periods]
If $f : \R \to \R$ is continuous and has a period-3 point, then $f$ has
periodic points of every period $n \in \N^*$.
\end{corollary}

\begin{remark}
Sharkovskii's ordering applies only in dimension one. In dimension two, a
diffeomorphism can have a period-3 point without having a fixed point (e.g.,
rotation by $2\pi/3$).
\end{remark}

\begin{lemma}[Intermediate value lemma]
\label{lem:covering}
Let $I, J$ be closed intervals and $f$ continuous. If $f(I) \supset J$, then
there exists a sub-interval $K \subset I$ with $f(K) = J$. In particular, if
$f(I) \supset I$, then $f$ has a fixed point in $I$.
\end{lemma}

\section{Symbolic Dynamics}

\begin{definition}[Sequence space]
\index{symbolic dynamics}
Let $\mathcal{A} = \{0, 1, \ldots, k-1\}$ be an alphabet of $k$ symbols.
The \textbf{sequence space} is
$\Sigma_k = \mathcal{A}^\N$, with metric
$d(s, t) = \sum_{i=0}^{\infty}\abs{s_i - t_i}/k^i$.
\end{definition}

\begin{definition}[Shift map]
\index{shift map}
The \textbf{shift} $\sigma : \Sigma_k \to \Sigma_k$ is
$\sigma(s_0, s_1, s_2, \ldots) = (s_1, s_2, s_3, \ldots)$.
\end{definition}

\begin{theorem}[Properties of the shift]
The shift $\sigma$ on $\Sigma_k$ satisfies:
\begin{enumerate}
    \item $\sigma$ is continuous and surjective (but not injective).
    \item Period-$n$ points are exactly the $n$-periodic sequences; there are
    $k^n$ of them.
    \item Periodic points are dense in $\Sigma_k$.
    \item $\sigma$ is topologically transitive.
    \item $\sigma$ is Devaney chaotic.
\end{enumerate}
\end{theorem}

\begin{definition}[Subshift of finite type]
\index{subshift of finite type}
Given a transition matrix $A \in \{0,1\}^{k \times k}$, the
\textbf{subshift of finite type} is
$\Sigma_A = \{s \in \Sigma_k : A_{s_i, s_{i+1}} = 1 \text{ for all } i\}$,
with $\sigma_A = \sigma|_{\Sigma_A}$.
\end{definition}

\begin{proposition}[Itinerary coding]
Let $f : [0,1] \to [0,1]$ be unimodal with critical point $c$. The
\textbf{itinerary} of $x$ is $s(x) = (s_0, s_1, \ldots)$ where $s_n = 0$
if $f^n(x) < c$ and $s_n = 1$ if $f^n(x) > c$. The map $x \mapsto s(x)$
semi-conjugates $f$ to the shift on a subspace of $\Sigma_2$.
\end{proposition}

\section{Topological Entropy}

\begin{definition}[Topological entropy]
\index{topological entropy}
Let $f : X \to X$ be continuous on a compact metric space. The
\textbf{topological entropy} is
\[
    h_{\mathrm{top}}(f) = \lim_{n \to \infty} \frac{1}{n}\ln s(n, \epsilon, f),
\]
where $s(n, \epsilon, f)$ is the maximum number of $(n, \epsilon)$-separated
points. The limit is independent of $\epsilon > 0$ sufficiently small.
\end{definition}

\begin{theorem}[Entropy of the shift]
The topological entropy of $\sigma$ on $\Sigma_k$ is
$h_{\mathrm{top}}(\sigma) = \ln k$. For the subshift of finite type
$\sigma_A$, $h_{\mathrm{top}}(\sigma_A) = \ln\rho(A)$, where $\rho(A)$ is
the spectral radius of $A$.
\end{theorem}

\begin{proposition}[Entropy of the logistic map]
For $f_4(x) = 4x(1-x)$, the topological entropy is
$h_{\mathrm{top}}(f_4) = \ln 2$. More generally, for the tent map of slope
$s$: $h_{\mathrm{top}} = \max(0, \ln s)$.
\end{proposition}

\begin{theorem}[Variational principle]
\index{variational principle}
For any continuous map $f$ on a compact metric space:
\[
    h_{\mathrm{top}}(f) = \sup_{\mu \in \mathcal{M}_f} h_\mu(f),
\]
where the supremum is over all $f$-invariant probability measures and
$h_\mu(f)$ is the metric entropy.
\end{theorem}

\begin{attention}[title=\textbf{Entropy and complexity}]
$h_{\mathrm{top}} > 0$ means the number of distinguishable orbits grows
exponentially with time. This is a more robust indicator of chaos than the
maximal Lyapunov exponent, as it is a topological invariant (independent of
the metric and measure).
\end{attention}

\section{Advanced Examples}

\begin{example}[Bernoulli shift]
The sawtooth map $B_k(x) = kx \pmod{1}$ on $[0,1)$ is exactly conjugate to
the shift on $\Sigma_k$ via the base-$k$ expansion. Its entropy is $\ln k$
and its Lyapunov exponent is $\ln k$.
\end{example}

\begin{example}[Quadratic Julia sets]
For $f_c(z) = z^2 + c$ with $c \in \C$, the Julia set
$J_c = \partial\{z : f_c^n(z) \not\to \infty\}$ is either connected (if $c$
is in the Mandelbrot set) or totally disconnected (a Cantor set). The
dynamics on $J_c$ displays rich chaotic phenomena.
\end{example}

\section{Exercises}

\begin{exercise}[Logistic bifurcations]
For the logistic map $f_r(x) = rx(1-x)$, analytically find the value of $r$
at which the period-2 cycle appears and the values of the cycle points.
\end{exercise}

\begin{exercise}[Sharkovskii]
Construct a continuous map $f : [0,1] \to [0,1]$ having a period-6 point but
no period-3 point. Show that it must have period-2 points.
\end{exercise}

\begin{exercise}[Symbolic dynamics]
Show that the shift on $\Sigma_2$ is topologically mixing: for all open sets
$U, V$, there exists $N$ with $\sigma^n(U) \cap V \neq \emptyset$ for all
$n \geq N$.
\end{exercise}

\begin{exercise}[Entropy]
Compute the topological entropy of the subshift of finite type defined by
$A = \begin{pmatrix} 1 & 1 \\ 1 & 0 \end{pmatrix}$ (Fibonacci shift). Show
that $h_{\mathrm{top}} = \ln\varphi$, where $\varphi = (1+\sqrt{5})/2$.
\end{exercise}

\begin{exercise}[Bifurcation diagram]
Write a program plotting the bifurcation diagram of the logistic map for
$r \in [2.5, 4]$. Visually identify the period-doubling cascade and periodic
windows (period 3, 5, 6, etc.).
\end{exercise}

\begin{exercise}[Kneading theory]
Determine the kneading itinerary of the critical point of $f_r$ for $r = 3.8$
and $r = 4$. Deduce which periods are forced.
\end{exercise}
